\chapter{Feasibility Analysis}
\section{Technical Feasibility}
The proposed system is technically feasible because it uses widely available and proven technologies such as the STM32F401 microcontroller, MEMS microphone, LM324 operational amplifiers, and LM386 audio amplifier. The analog circuits handle most of the noise suppression, while a lightweight TinyML model performs transient noise detection, reducing computational requirements and enabling real-time operation with low latency.
\section{Economic Feasibility}
The proposed Real-Time Assistive Hearing System is economically feasible because it uses readily available and low-cost electronic components that can be sourced locally. The hybrid analog-digital architecture reduces the need for expensive signal-processing hardware by performing most noise suppression through analog circuits. Additionally, the project relies on freely available software tools such as STM32CubeIDE, STM32Cube.AI, and TensorFlow, making development accessible and cost-effective. Its modular design also simplifies maintenance and future upgrades, further enhancing its practicality and sustainability.
\section{Operational Feasibility}
The system is practical and easy to operate for its intended users, including industrial workers, security personnel, and individuals with mild conductive hearing loss. By using bone-conduction output, the device allows users to hear amplified speech while keeping their ears open to environmental sounds, improving both communication and situational awareness.
\section{Schedule Feasibility}
The project can be completed within the duration of a typical engineering project cycle. Activities such as circuit design, hardware assembly, TinyML model training, firmware development, testing, and documentation can be carried out systematically within the available academic timeline.
\section{Resource Feasibility}
All required resources, including electronic components, development software, and laboratory equipment such as oscilloscopes and soldering stations, are readily accessible. The project team also possesses the necessary knowledge of electronics, embedded systems, and machine learning to successfully implement and test the system.
\section{Overall Feasibility}
Overall, the proposed assistive hearing system is highly feasible from technical, economic, operational, schedule, and resource perspectives. Its low-cost design, use of readily available components, and practical real-world applications make it a suitable and achievable engineering project.
